\documentclass[10pt]{article}

\usepackage[margin=0.9in,top=0.8in,bottom=0.8in]{geometry}
\usepackage{amsmath,amssymb}
\usepackage{listings}
\usepackage{xcolor}
\usepackage{hyperref}
\usepackage{enumitem}
\usepackage{booktabs}

\definecolor{codebg}{rgb}{0.96,0.96,0.96}
\definecolor{codekw}{rgb}{0.0,0.0,0.6}
\definecolor{codecomment}{rgb}{0.4,0.4,0.4}
\definecolor{codestring}{rgb}{0.58,0.0,0.0}

\lstdefinelanguage{Julia}{
  morekeywords={function,end,for,while,if,else,elseif,begin,return,struct,
    where,do,in,let,const,mutable,abstract,type,import,using,module,
    export,true,false,nothing,Inf,NaN},
  sensitive=true,
  morecomment=[l]{\#},
  morestring=[b]",
}

\lstset{
  backgroundcolor=\color{codebg},
  basicstyle=\ttfamily\footnotesize,
  keywordstyle=\color{codekw}\bfseries,
  commentstyle=\color{codecomment}\itshape,
  stringstyle=\color{codestring},
  numbers=none,
  breaklines=true,
  showstringspaces=false,
  frame=single,
  rulecolor=\color{gray!40},
  xleftmargin=1em,
  xrightmargin=1em,
  aboveskip=0.4em,
  belowskip=0.1em,
  tabsize=2,
}

\usepackage{titlesec}
\titlespacing*{\section}{0pt}{1.1em}{0.5em}
\titlespacing*{\subsection}{0pt}{0.9em}{0.4em}
\setlength{\parskip}{0.15em}

\newcommand{\ESS}{\mathrm{ESS}}
\newcommand{\Lhat}{\hat L}

\title{Two Julia Implementations of the Trigger/Target Particle Filter:\\
A Comparison}
\author{Prepared for the PartiallyObservedMarkovProcesses.jl project}
\date{September 8, 2026}

\begin{document}
\maketitle

\begin{abstract}
Both repositories implement the same generalized sequential Monte Carlo
(particle filter) recursion for a partially observed Markov process:
resampling is triggered by a threshold on the effective sample size, and,
when it occurs, is tempered by a power $\beta$ that interpolates between
equally weighted resampling and an untouched weighted particle
representation. This note compares Aaron King's Julia translation
(\texttt{PartiallyObservedMarkovProcesses.jl\_aarons/src/pfilter.jl}, at
commit \texttt{1610d34}) against Debsurya's Julia translation
(\texttt{PartiallyObservedMarkovProcesses.jl/src/pfilter.jl}), which share
the same trigger/target interface and the same recursive structure. The
principal finding is that power-tempered resampling discards a normalizing
constant $C$ whenever $0<\beta<1$: dropping it, as Aaron's Julia translation
and the R \texttt{pomp} package's \texttt{wpfilter.c} both do, biases the
running log likelihood low by $\Theta(1/n)$ per resampling step, a claim
established here by an exact two-particle calculation and confirmed by an
end-to-end test built specifically to be sensitive to it; Debsurya's
translation restores $C$. The note also documents an indexing error, present
in the compared Aaron King Julia translation, that assigned the retained
weight by resampling slot rather than by selected ancestor. Aaron King has
since fixed this himself, independently, in commit \texttt{1d8dbc0}, in a
form equivalent to Debsurya's, so it is recorded here as a resolved
historical point --- and the comparison against it explicitly as a
comparison against the pre-\texttt{1d8dbc0} code --- rather than as a live
divergence. Two further differences are documented: an unequally weighted
terminal ancestry draw, which Aaron's code initiated uniformly rather than
proportional to the terminal weights, and the presence, in Debsurya's
translation only, of a per-step resampling diagnostic. It closes
with the resulting unification of \texttt{wpfilter} as a thin wrapper around
\texttt{pfilter}.
\end{abstract}
\vspace{0.6em}
\noindent\textbf{Status, 21 September 2026.}\; Every finding in this note has
since been fixed upstream, following correspondence with its author: the
normalizing constant in \texttt{3dac3b5} (``ensure unbiasedness of pfilter'',
16 September) and the terminal ancestry draw in \texttt{24febe8}
(9 September). The upstream correction is arithmetically identical to the one
described here, written against a unit-sum rather than a unit-mean weight
convention --- $\log(s\cdot du)$ with $s=\sum_j w_{A_j}^\beta$ and $du=S/n$,
in place of $\log(mS/n)$; the two agree to $1.4\times10^{-16}$ on the
two-particle case below. What follows is therefore a record of the argument
and of why the obvious tests missed it, not a description of a current
divergence.


\section{The trigger/target resampling scheme}

Both implementations run the same filtering recursion. At each observation
time, particles $x_1,\dots,x_n$ carry linear weights $w_1,\dots,w_n$,
normalized so that $\bar w = n^{-1}\sum_j w_j = 1$. The effective sample size
\[
  \ESS \;=\; \frac{\left(\sum_j w_j\right)^2}{\sum_j w_j^2}
\]
is compared against a threshold set by the \texttt{trigger} parameter,
$\mathrm{trigger}\in[0,1]$: resampling occurs whenever $\ESS \le
\mathrm{trigger}\cdot n$. \texttt{trigger}$=1$ resamples at every
observation time (the classical bootstrap filter); \texttt{trigger}$=0$
never resamples, and the incremental weights are simply carried forward,
multiplied into the running weight vector at the next observation time.

When resampling does occur, it is tempered by the \texttt{target} parameter
$\beta\in[0,1]$: particle $j$ is selected as an ancestor with probability
proportional to $w_j^{1-\beta}$, and each selected particle retains, going
forward, a weight proportional to $w_j^{\beta}$ (the weight of the
particular ancestor selected), renormalized to unit mean across the new
particle set. This is a genuine one-parameter family: at $\beta=0$, ancestors
are drawn proportional to the full weight $w_j$ and every survivor is
reweighted to $1$ — the standard bootstrap resampling step. At $\beta=1$,
ancestors are drawn without regard to weight ($w_j^0=1$ for all $j$), and
each survivor retains its ancestor's original weight $w_j^1=w_j$ — since
uniform systematic selection of $n$ ancestors from $n$ slots reproduces the
identity permutation up to a vanishing boundary effect, this carries the
full weighted particle representation forward essentially unchanged. Both
repositories implement this design; the divergences below concern its
realization in code, not the design itself.

\section{A positional-indexing error in weight retention, since fixed upstream}
\label{sec:positional}

This section compares against Aaron's Julia translation as it stood at
commit \texttt{1610d34}, the version installed locally when this comparison
was prepared; it is not a live difference between the two repositories
today; see the note at the end of the section. The underlying reasoning is
kept because it is used again, unchanged, in \S\ref{sec:normconst} below.

Once ancestor indices $p(1),\dots,p(n)$ have been drawn from the cumulative
sums of $w_j^{1-\beta}$, the retained weight $w_j^\beta$ must travel
\emph{with the particle it came from}: since $w_j = w_j^{1-\beta}\cdot
w_j^\beta$ and the first factor has already been consumed by the selection
probability, the second factor is what remains to be assigned to the
particle now occupying position $k=p^{-1}(j)$ in the resampled array.

Aaron's Julia \texttt{systematic\_resample!} instead raises the entire
weight vector to the power $\beta$ \emph{in place, in its original
positional order}, after the ancestor indices have already been drawn:

\begin{lstlisting}[language=Julia, caption={Aaron's \texttt{systematic\_resample!} (PartiallyObservedMarkovProcesses.jl\_aarons/src/pfilter.jl, lines 336--363)}]
systematic_resample!(
    p::AbstractArray{I,1},
    w::AbstractArray{W,1},
    ucum::AbstractArray{W,1},
    beta::Float64,
) where {I,W} = begin
    s::W = 0
    alpha = 1-beta
    @inbounds for j in eachindex(w)
        s += w[j]^alpha
        ucum[j] = s
    end
    n::I = length(w)
    i::I = 1
    du::W = s/n
    u::W = -du*rand(W)
    @inbounds for j in eachindex(p)
        u += du
        while (u > ucum[i] && i < n)
            i += 1
        end
        p[j] = i
    end
    w .^= beta
    w ./= mean(w)
    nothing
end
\end{lstlisting}

The particle \emph{states} are copied correctly downstream
(\texttt{xf .= xp[p]}, so state $k$ correctly becomes the state of ancestor
$p(k)$), but the weight vector \texttt{w} is exponentiated positionally:
the new weight at position $k$ is $w_k^\beta$ — the $\beta$-power of
whatever particle \emph{used to occupy} position $k$ — not $w_{p(k)}^\beta$,
the $\beta$-power of the ancestor actually selected. Whenever the
resampling permutation $p$ moves mass between positions, which is the norm
rather than the exception, the state and its retained weight come from two
different particles.

Aaron's R implementation of the same routine (the C code behind
\texttt{wpfilter} in the R \texttt{pomp} package) does not have this defect:
it copies the retained weight of the selected ancestor directly, so the
weight is indexed by the ancestor, not by position:

\begin{lstlisting}[language=C, caption={R \texttt{wpfilter.c}, inside the ancestor-selection loop, as reported by Debsurya}]
wt[k] = ws[sp];    /* ws = carried weights, sp = selected ancestor index */
\end{lstlisting}

Debsurya's Julia translation takes the ancestor-indexed form, following the
R form rather than the positional form above:

\begin{lstlisting}[language=Julia, caption={Debsurya's \texttt{systematic\_resample!} (PartiallyObservedMarkovProcesses.jl/src/pfilter.jl, lines 457--491), retained-weight step}]
@inbounds for j in eachindex(p)
    ucum[j] = w[p[j]]^beta     # retained weight of the selected ancestor
end
@inbounds for j in eachindex(w)
    w[j] = ucum[j]
end
m::W = mean(w)
w ./= m
log(m*s/n)
\end{lstlisting}

Here \texttt{w[p[j]]} is evaluated \emph{before} \texttt{w} is overwritten,
using \texttt{ucum} as scratch space, so each new weight is correctly
indexed by the selected ancestor. At $\beta=0$ the two forms coincide
trivially (every weight is $1$ regardless of index), which is presumably
why the error had not surfaced in the classical ($\mathrm{target}=0$)
filter; it is only visible for $\beta>0$.

\paragraph{This is no longer a difference.} Aaron King fixed the positional
form himself, independently, in commit \texttt{1d8dbc0} (``fix bug in
weighted particle filter,'' 2026-09-08), which postdates \texttt{1610d34}
and is not an ancestor of the local clone compared above:
\begin{lstlisting}[language=Julia, caption={Aaron's \texttt{systematic\_resample!} after commit \texttt{1d8dbc0}, retained-weight step}]
@inbounds for j ∈ eachindex(p)
    ucum[j] = w[p[j]]
end
@inbounds for j ∈ eachindex(w)
    w[j] = ucum[j]^β
end
\end{lstlisting}
which is equivalent to Debsurya's form: both index the retained weight by
the selected ancestor $p(k)$ rather than by resampling slot $k$. The finding
above is accordingly recorded as a resolved historical point, not a present
divergence between the two implementations. It is kept in this note in full
because the same reasoning about which quantity must travel with the
selected ancestor, rather than with the slot it lands in, is exactly what
\S\ref{sec:normconst} needs to derive the normalizing constant discarded at
partial resampling.

\section{The normalizing constant discarded at partial resampling}
\label{sec:normconst}

Write $w_1,\dots,w_n$ for the unit-mean weights immediately prior to
resampling, and $S = \sum_j w_j^{1-\beta}$ for the normalizing sum of the
selection probabilities $q_j = w_j^{1-\beta}/S$. Because ancestors are drawn
from $q$ rather than from the weights themselves, the sampling measure has
changed, and the properly weighted representation assigns the particle
selected as ancestor $A_k$ at position $k$ the importance weight
\[
  R_k \;=\; \frac{w_{A_k}}{n\,q_{A_k}} \;=\; \frac{S}{n}\,w_{A_k}^{\beta},
\]
whose sample mean over the $n$ selections,
\[
  C \;=\; \frac1n\sum_{k=1}^n R_k \;=\; \frac{S}{n}\,m,
  \qquad m = \frac1n\sum_{k=1}^n w_{A_k}^{\beta},
\]
is \emph{the resampling step's own contribution to the normalizing constant}
of the unnormalized (Feynman--Kac) measure the filter estimates. Both
implementations instead store the retained weights renormalized to unit
mean, $V_k = R_k/C$ — necessary since every other routine in the filter (the
effective-sample-size and conditional-log-likelihood computation in
particular) assumes a unit-mean weight vector on entry — and neither
multiplies $C$ back into the likelihood accumulator. Aaron's Julia routine
returns \texttt{nothing}, its caller discarding it:
\begin{lstlisting}[language=Julia, caption={Aaron's \texttt{pfilt\_step\_comps!} (PartiallyObservedMarkovProcesses.jl\_aarons/src/pfilter.jl, line 228)}]
systematic_resample!(p, w, work, target)   # return value discarded
\end{lstlisting}
and R \texttt{pomp}'s \texttt{wpfilter.c} discards it too, by a different
route: it leaves the carried weights unnormalized and computes each
conditional log likelihood as a ratio of successive \emph{total} masses
(\texttt{loglik = maxlogw + log(w) - log(oldw)}, line 110), so that the mass
change across a resampling step is simply never entered into the likelihood
at all — it becomes part of the baseline the next step's ratio is taken
against, with no correction of any kind at the step where it arose. $C$ is
dropped, not restored, in both of Aaron King's implementations, and, until
now, in this repository as well.

\subsection*{Why dropping $C$ is a real bias}

$C$ is conditionally mean one: $\mathbb E[m\mid w] = \sum_j q_j\,w_j^\beta =
(\sum_j w_j)/S = n/S$, so $\mathbb E[C\mid w]=1$ exactly. It is tempting to
conclude from this alone that dropping $C$ costs only variance, since a
factor with conditional expectation one, multiplied into an
otherwise-unbiased estimator, cannot change that estimator's expectation.
That inference needs independence, or at least the absence of the wrong kind
of correlation, and $C$ does not have it: $C$ is a function of which
ancestors $A_{1:n}$ were actually selected, and those same ancestors go on
to generate \emph{every future contribution to the likelihood}. Knowing
$\mathbb E[C]=1$ does not give $\mathbb E[C\,\ell]=\mathbb E[\ell]$ once $C$
and $\ell$ are dependent through a shared cause — here, the resampled
particle cloud itself. Dropping $C$ biases the likelihood estimate downward
by $\Theta(1/n)$ per resampling step. $C\equiv1$ identically at $\beta=0$
(where $S=\sum_j w_j=n$ and every retained weight is $1$, so $m=1$) and at
$\beta=1$ (where $S=n$ trivially and, since uniform systematic selection of
$n$ ancestors from $n$ slots reproduces the identity permutation, $m =
n^{-1}\sum_j w_j = 1$ as well), so the ordinary bootstrap filter is entirely
untouched: this is strictly a partially-retained-weight ($0<\beta<1$)
phenomenon.

\subsection*{An exact counterexample}

Take $n=2$ particles with unit-mean weights $w=(1.8,0.2)$ and $\beta=1/2$.
Then $q=(0.75,0.25)$, and two-particle systematic resampling gives ancestry
$(1,1)$ or $(1,2)$, each with probability exactly $1/2$, carrying $C=1.2$
and $C=0.8$ respectively. Suppose the future contribution from the two
particles is $g=(1,0)$; the correctly weighted predictive quantity is
$\frac1n\sum_i w_i g_i = 0.9$. Dropping $C$ gives
$\frac12(1)+\frac12(0.75)=0.875$; retaining it gives
$\frac12(1.2\cdot1)+\frac12(0.8\cdot0.75)=0.9$, exactly. This has been
verified numerically to six decimal places over $4{,}000{,}000$ replicates.

\subsection*{An end-to-end test, and why an earlier check missed this}

\texttt{test/iid.jl} contains two test sets. The first uses a latent process
with no memory: $X_t$ is drawn afresh, independently of $X_{t-1}$, at every
step. There the one-step predictive density is the same for every particle
regardless of which ancestors survive resampling, so the tempered
resampling step is exactly unbiased whatever it does with $C$ — this model
is blind to the effect by construction, at any number of replicates. The
second test uses a frozen binary latent state, $X\sim\mathrm{Bernoulli}(1/2)$
with $X_t=X$ for all $t$, together with measurement potentials
$g_1(1)=1.8,\,g_1(0)=0.2$ and $g_2(1)=1.0,\,g_2(0)=0.0$; the exact likelihood
is $Z=\mathbb E[g_1(X)g_2(X)]=0.9$. At $n=2$ and $\beta=1/2$, dropping $C$
gives $\mathbb E[\hat Z]\approx0.8879$ against $Z=0.9$ — a $27$-standard-error
discrepancy over $200{,}000$ replicates — that vanishes when $C$ is restored.

An earlier check, run instead on the linear-Gaussian reduction of the
Gompertz model against its closed-form Kalman likelihood over $400{,}000$
paired replicates at $\beta\in\{0.25,0.5,0.75\}$ and $n\in\{8,16,64\}$, found
no significant difference between the two conventions and was taken,
wrongly, as evidence that $C$ needed no restoring. That check had
no power: the bias is $\Theta(1/n)$ per resampling step and vanishes
identically whenever the one-step predictive density is constant across the
cloud, which is very nearly true of a smooth, well-mixing latent process
such as log-Gompertz. The design could resolve differences of roughly
$0.02\%$ of $\hat L$ against a bias under $0.15\%$ — an underpowered test, not
a clean null result. A null Monte Carlo comparison is evidence of nothing
unless the test model can exhibit the effect in question and the design has
power against it; the frozen-state model above was built specifically to
have both properties.

\subsection*{The fix}

Debsurya's \texttt{systematic\_resample!} now returns $\log C$ instead of
\texttt{nothing}, and the caller accumulates it into the running log
likelihood:
\begin{lstlisting}[language=Julia, caption={PartiallyObservedMarkovProcesses.jl/src/pfilter.jl, lines 481--490 and 315}]
    stot::W = s   # S = sum_j w_j^(1-beta)
    ...
    m::W = mean(w)
    w ./= m # subsequent steps rely on the weights having unit mean
    log(m*stot/n)
\end{lstlisting}
\begin{lstlisting}[language=Julia, caption={The caller, \texttt{pfilt\_step\_comps!}, line 315}]
        logLik[] += systematic_resample!(p, w, work, target)
\end{lstlisting}
Aaron's Julia translation and R \texttt{pomp}'s \texttt{wpfilter.c} still
drop $C$, unchanged from the description above.

\section{Terminal ancestry draw under unequal terminal weights}
\label{sec:ancestry}

Whenever $\beta>0$, or a step is skipped by \texttt{trigger}, the terminal
particle cloud carries unequal weights: the filtering distribution's
information sits partly in the weights, not solely in the particle
multiplicities. Recovering a single properly weighted draw of the filtered
trajectory therefore requires initiating the traced lineage at the terminal
time with probability proportional to the terminal weights, then tracing
ancestry backward through the recorded permutations.

Aaron's \texttt{trace\_ancestry!}, at the revision compared against here,
initiated the draw uniformly, regardless of the terminal weights (fixed
upstream in \texttt{24febe8}):
\begin{lstlisting}[language=Julia, caption={Aaron's \texttt{trace\_ancestry!} (PartiallyObservedMarkovProcesses.jl\_aarons/src/pfilter.jl, line 397)}]
j::I = rand(axes(perm,2))     # uniform draw, regardless of terminal weights
\end{lstlisting}
A marginal comment in that file (\texttt{\# FIXME: check that the first
index in the ancestry is correct}) suggests this was already a point of
concern. Unlike the positional-indexing error of \S\ref{sec:positional},
this one survives commit \texttt{1d8dbc0}: the terminal draw there (line
408) is still the uniform \texttt{j = rand(axes(perm,2))}, and the FIXME
comment (now at line 97) is still unresolved. This is accordingly a live
difference between the two implementations, not a historical one.

A minimal counterexample makes the resulting bias concrete. Suppose two
terminal particles carry weights $W=(0.9,\,0.1)$ — an entirely ordinary
outcome once $\beta>0$ or a step goes unresampled. The correct marginal
draw of the terminal state assigns particle $1$ probability $0.9$ and
particle $2$ probability $0.1$, so that a functional such as the traced
trajectory's terminal value has, e.g., $\mathbb P(X_T = x^{(1)}) = 0.9$.
Aaron's uniform draw instead assigns each particle probability $0.5$
regardless of $W$, systematically over-weighting the low-probability
particle and under-weighting the high-probability one; the traced
trajectory is properly weighted for the filtering distribution only by
coincidence, when the terminal weights happen to be equal.

Debsurya's \texttt{trace\_ancestry!} initiates the draw proportional to the
terminal weights before tracing backward:
\begin{lstlisting}[language=Julia, caption={Debsurya's weighted \texttt{trace\_ancestry!} (PartiallyObservedMarkovProcesses.jl/src/pfilter.jl, lines 518--554), initiation step}]
s::W = 0
for k in eachindex(w)
    v::W = w[k]
    s += v
    work[k] = s
end
u::W = s*rand(LogLik)
jj::I = 1
while (u > work[jj] && jj < n)
    jj += 1
end
j = jj
# ... trace backward from j exactly as before
\end{lstlisting}
with a fallback to the uniform draw only in the degenerate case where every
terminal weight underflows to zero. This routine was exercised in the
model-based validation against R \texttt{pomp} and the closed-form Gompertz
likelihood described elsewhere, and directly addresses the concern flagged
in Aaron's comment.

\section{A per-step resampling diagnostic}
\label{sec:resampled}

Debsurya's \texttt{PfilterdPompObject} carries an additional field,
\texttt{resampled::Array\{Bool,1\}}, recording at each observation time
whether systematic resampling actually occurred there (as opposed to being
skipped because $\ESS$ stayed above the \texttt{trigger} threshold), together
with an accessor:
\begin{lstlisting}[language=Julia, caption={PartiallyObservedMarkovProcesses.jl/src/pfilter.jl, lines 29--35}]
"""
    resampled(object)

`resampled` extracts the vector of booleans recording, for each
observation time, whether systematic resampling occurred.
"""
resampled(object::PfilterdPompObject) = object.resampled
\end{lstlisting}
This is set alongside the ancestry permutation in \texttt{pfilt\_step\_comps!}
(\texttt{resamp[] = true} when resampling occurs, \texttt{false} otherwise)
and gives a cheap, per-step diagnostic of weight degeneracy over the course
of a filtering run, independent of and complementary to the effective
sample size trace. Aaron's \texttt{PfilterdPompObject} has no corresponding
field.

\section{Unification: \texttt{wpfilter} as a thin wrapper}
\label{sec:unification}

In both repositories, the classical bootstrap filter
($\mathrm{trigger}=1,\ \mathrm{target}=0$) and the general trigger/target
filter were, at one point, separately implemented engines: a fast path with
no weight-carrying overhead for the classical case, and a general path that
threads a carried weight vector \texttt{wcarry} through the recursion (folded
into the incremental log weights at \texttt{compute\_ess\_logLik!}, and
consulted at each resampling step). Debsurya's \texttt{pfilter} still
dispatches internally between these two engines on the basis of whether
\texttt{trigger==1 \&\& target==0}, but both are reached through a single
public entry point, and \texttt{wpfilter} — the name inherited from the R
\texttt{pomp} package's separate weighted-filter routine — is now nothing
more than a same-signature call-through to \texttt{pfilter}:
\begin{lstlisting}[language=Julia, caption={PartiallyObservedMarkovProcesses.jl/src/wpfilter.jl, in full}]
wpfilter(
    object::ValidPompData;
    trigger::Real = 1,
    target::Real = 0,
    kwargs...,
) = pfilter(object; trigger, target, kwargs...)

wpfilter(
    object::PfilterdPompObject;
    Np::Integer = object.Np,
    trigger::Real = object.trigger,
    target::Real = object.target,
    kwargs...,
) = pfilter(object; Np, trigger, target, kwargs...)

wpfilter(_...) = error("Incorrect call to `wpfilter`.")
\end{lstlisting}
There is exactly one particle-filtering implementation to validate,
maintain, and reason about; \texttt{wpfilter} survives purely as an
API-compatible alias, with defaults ($\mathrm{trigger}=1$,
$\mathrm{target}=0$) that in fact route it to the same classical engine as a
bare call to \texttt{pfilter}. This unification is also what makes it
possible for \texttt{mif2}, the iterated-filtering routine, to run its
inner filtering step under $\mathrm{target}>0$: since there is only one
filtering routine underneath both names, a weighted filtering step is
automatically available wherever a classical one was, which is not presently
true of the R implementation, whose weighted filter does not resample
parameters.

\section{Summary}

\begin{center}
\begin{tabular}{@{}p{0.32\textwidth}p{0.32\textwidth}p{0.32\textwidth}@{}}
\toprule
\textbf{Issue} & \textbf{Aaron's Julia} & \textbf{Debsurya's Julia} \\
\midrule
Retained weight indexing (\S\ref{sec:positional})
  & ancestor-indexed \newline\emph{(positional at \texttt{1610d34}; fixed independently at \texttt{1d8dbc0})}
  & ancestor-indexed: $w_{p(k)}^\beta$ \newline\emph{(no longer a difference)} \\[0.4em]
Normalizing constant $C$ at $0<\beta<1$ (\S\ref{sec:normconst})
  & dropped at \texttt{1d8dbc0} \newline\emph{(biases $\hat L$ by $\Theta(1/n)$ per step; fixed upstream in \texttt{3dac3b5})}
  & restored: \texttt{logLik[] += log C} \\[0.4em]
Terminal ancestry draw (\S\ref{sec:ancestry})
  & uniform at \texttt{1d8dbc0} \newline\emph{(fixed upstream in \texttt{24febe8})}
  & proportional to terminal weight \\[0.4em]
Per-step resampling flag (\S\ref{sec:resampled})
  & absent
  & \texttt{resampled(object)} \\[0.4em]
\texttt{wpfilter} (\S\ref{sec:unification})
  & separate historical engine
  & thin wrapper over \texttt{pfilter} \\
\bottomrule
\end{tabular}
\end{center}

\noindent Of the five rows, only the third and fourth are differences of the
comparison sort this note set out to find: for $\mathrm{target}>0$, the
terminal ancestry draw changes the distribution represented by the stored
filtered trajectory, and the resampling flag is simply present in one
implementation and not the other. Neither disturbs the marginal
log-likelihood, so neither would have been caught by comparing likelihood
estimates; they were found by reading the code. The first row was of the
same kind — it would have changed which particle's information survives
resampling for $\mathrm{target}>0$ — but Aaron King fixed it independently
before this note was finished, so it is recorded as resolved rather than
open. The second row is different in kind from all the others: it is not a
difference of code structure but a shared numerical defect, present in
Aaron's Julia translation and in R \texttt{pomp} alike, that this
repository has now corrected. Unlike the first and third rows, it
\emph{does} disturb the marginal log-likelihood, by $\Theta(1/n)$ per
resampling step at every $0<\beta<1$; the reason it went undetected for as
long as it did is that the model used to check the filter's likelihood
against an exact reference happened to be almost exactly the case,
described in \S\ref{sec:normconst}, in which the effect vanishes.

\end{document}
